\hypertarget{front-matter}{%
\label{front-matter}}

\hypertarget{abstract}{%
\chapter*{Abstract}
\addcontentsline{toc}{chapter}{Abstract}
\markboth{Abstract}{Abstract}\label{abstract}}

This book sets out to derive Graham Priest's non-classical and dialetheic
logics from a general theory of constrained representation, rather than
presenting classical logic as the default and non-classical systems as
deviations from it. It does not deliver the result it initially expected.
Early planning aimed at a Glut--Obstruction Correspondence: some precise
identification of Priest's contradictions with a geometric obstruction to
assembling local evidence into a coherent whole. The completed manuscript
establishes close to the opposite, and the result generalizes past the
one theorem it was first stated for: \textbf{contradiction and obstruction are
independent facts about the same data.} Chapter 22's Mutual
Non-Determination Proposition is this claim's sharpest instance --- whether
a proposition is contradictory and whether its supporting evidence
coheres do not determine one another, and an explicit six-cell
construction proves it rather than merely failing to disprove it --- but
the same pattern recurs throughout the book: apparently related
structures are repeatedly shown to carry independent coordinates rather
than collapsing into one. Three further correspondences investigated
with equal seriousness --- non-normal worlds as refusal, bilattice merge as
ordinary sheaf gluing, glut as a nonzero cohomology class --- were rejected
outright, at three different structural levels (modal, operational,
geometric), after being plausible enough to have been provisionally
accepted earlier in the project. The book's actual thesis, assembled from
these results rather than imposed on them in advance, is that
non-classical logics need not be understood as competing relaxations of
one classical standard; they can be understood as responses to different
failures along structural axes that turn out, on inspection, to be
substantially independent of one another.

\hypertarget{introduction}{%
\chapter*{Introduction}
\addcontentsline{toc}{chapter}{Introduction}
\markboth{Introduction}{Introduction}\label{introduction}}

\hypertarget{what-this-book-does-not-do}{%
\section{What this book does not do}\label{what-this-book-does-not-do}}

It does not treat classical logic as the ground state from which
paraconsistent, intuitionistic, relevant, many-valued, and modal systems
are departures requiring justification. Chapter 1 states the alternative
directly: logical consequence is preservation under admissible
transformation, \(x \leadsto_R y\), with classical logic sitting at one
specific point in the space of possible regimes \(R\) --- the point
combining maximal strictness on valuation (no gaps, no gluts) with
unrestricted explosive propagation from contradiction under the relevant
closure principles. Every other logic covered in this book is a
different point in the same space --- an alternative regime whose adequacy
depends on the constraints governing the representational task at hand,
not an equally good option regardless of task, and not a weakening of
the classical one. Admissibility is precisely what lets the book
discriminate among regimes; nothing here is an argument for treating all
regimes as interchangeable.

By the book's own account, this way of stating things is only earned
because three plausible correspondences failed to survive their
chapters: non-normality is not refusal, merge is not gluing, and glut is
not obstruction. A reader should not expect a book that translates every
Priest concept into this framework's vocabulary and declares victory ---
several translations were attempted and explicitly rejected, and that
outcome is treated as evidence for the book's thesis rather than as
incidental failure.

\hypertarget{the-question-the-book-actually-asks}{%
\section{The question the book actually asks}\label{the-question-the-book-actually-asks}}

Stated once, in Chapter 1, and returned to throughout: \emph{what must remain
invariant when representations are combined, transported, refined,
contradicted, or extended?} This is deliberately not the question
``which logic is correct.'' It's a question about which structural
distinctions survive contact with an actual formal system, once that
system is built rather than merely proposed.

\hypertarget{separation-before-unification}{%
\section{Separation before unification}\label{separation-before-unification}}

The book's method, visible in retrospect across all six parts, is to
take something ordinarily treated as one notion and ask whether it's
actually several. Nearly every chapter that matters does this. Chapter 6
separates admissibility and accessibility: what states a regime permits,
versus which of those states a given world can modally reach. Chapter 7
separates inconsistency, impossibility, and non-normality: a state
containing a contradiction, a state failing to satisfy a governing
regime, and a state whose evaluation rules are themselves suspended are
three independent facts, not one phenomenon under three names. Chapters
9 and 10 separate positive and negative support: the two questions ``is
there evidence for'' and ``is there evidence against,'' whose four
possible joint answers generate Belnap and Dunn's four truth values as a
consequence rather than a stipulation. Chapter 13 separates support and
constructibility: whether evidence exists, and whether it takes the
specific form of an exhibitable witness, answer different questions even
when the same proposition is at issue. Chapters 15 and 20--21 separate
compatibility and dispersion: local evidence disagreeing outright, versus
locally negligible variation that nonetheless accumulates past a
threshold under coarse-graining, are different failure mechanisms with
different mathematics. Chapters 20--22 separate semantic glut from
geometric coherence defect: a proposition's contradictory support status,
and whether the evidence behind that status agrees with itself across
overlapping contexts, do not determine one another. Chapters 20--21
separate gluing, merge, and repair: reconstructing a section from a
family that already agrees, versus combining a family that doesn't agree
while retaining everything it supplied, are different operations that an
earlier draft of this project once conflated under one word.

None of these separations were chosen for elegance in advance. Several
were forced by a correspondence failing to hold up once its chapter was
actually drafted --- recorded, not smoothed over, in the correspondence
ledger this book's technical claims are checked against.

\hypertarget{the-hierarchy-this-produces}{%
\section{The hierarchy this produces}\label{the-hierarchy-this-produces}}

Six questions, asked of the same underlying material at increasing
levels of structure, form the closest thing this book has to a single
conceptual synopsis:

\[\text{support} \to \text{construction} \to \text{compatibility} \to \text{repair} \to \text{persistence} \to \text{continuation}.\]

Does evidence exist? Can it be exhibited? Can local descriptions coexist?
If not, can they be minimally reconciled without discarding either? Does
the resulting determination survive a change of scale? Does it remain
identifiable as the same thing through a change of state? These are
successive questions in the book's analysis, asked in this order because
each becomes askable once the last has been separated out --- not because
each presupposes or reduces to the one before it. Construction does not
intrinsically presuppose support, and compatibility does not
intrinsically presuppose construction; the book argues for their
independence rather than building a reduction chain. The six parts
largely track this order, though not in the sequence they were drafted.

\hypertarget{the-parameter-that-becomes-structure}{%
\section{The parameter that becomes structure}\label{the-parameter-that-becomes-structure}}

There's a second way to see the same six parts, sharper for being purely
technical: watch what happens to the single parameter \(R\) introduced in
Chapter 1.

\[R \to (R, \mathcal{R}) \to (R, \mathcal{R}, \sigma^+, \sigma^-) \to \text{local support/coherence} \to \text{repair and persistence} \to R_w\text{-dependent continuation}.\]

Chapter 1 treats admissibility as one relation governing consequence.
Chapter 6 finds that modal reasoning needs a second, independent
relation, \(\mathcal{R}\), accessibility. Chapters 9--10 find that a single
admissibility judgment isn't enough to generate FDE's four values without
splitting support into two independent coordinates, \(\sigma^+\) and
\(\sigma^-\). Part VI finds that those coordinates, once localized across a
cover, need both a support structure and a separate coherence structure
to say what a family of local reports actually agrees on. Chapter 23's
repair functional finds that action depends on both structures
independently, via a term the earlier two-term cost function had no room
for. And Chapter 7's \(R_w\) --- needed to place non-normal worlds correctly
--- turns out, once Part V asks about identity across time, to mean that
the admissibility regime itself can vary along a continuation, which is
what produces this book's one genuine open problem rather than a known
next step. The apparent single parameter of Chapter 1 does not stay a
parameter. It progressively unpacks into structure, and the book's final
open question is what happens once the last piece of that structure ---
\(R\) itself --- stops holding still.

\hypertarget{the-methodological-thesis}{%
\section{The methodological thesis}\label{the-methodological-thesis}}

\begin{quote}
Non-classical logics need not be understood as competing relaxations of
one classical standard. They can instead be understood as responses to
different failures or constraints along partially independent
structural axes.
\end{quote}

Paraconsistency, on this account, is what a regime looks like once it
stops requiring gluts to be classically impossible. Intuitionism is what
happens once construction is required in addition to support. Vagueness
is what a resolution-indexed system looks like once persistence under
coarse-graining, not mere agreement, becomes the relevant question. None
of these needs to be motivated as a retreat from classical logic's
demands; each answers a genuinely different question the classical
regime never had to ask because it collapsed all these axes into one.

\hypertarget{what-this-introduction-promises-and-no-more}{%
\section{What this introduction promises, and no more}\label{what-this-introduction-promises-and-no-more}}

Every technical claim referenced above is checked against the
correspondence ledger, which records status after drafting rather than
before --- including three correspondences the book investigated and
rejected, and one genuine open problem the completed manuscript produced
rather than resolved. This introduction does not promise more than that
ledger backs.

\hypertarget{terminology-table}{%
\chapter*{Terminology Table}
\addcontentsline{toc}{chapter}{Terminology Table}
\markboth{Terminology Table}{Terminology Table}\label{terminology-table}}

Pairs the book keeps distinct, and what conflating them would cost.

\begin{longtable}[]{@{}
  >{\raggedright\arraybackslash}p{(\columnwidth - 4\tabcolsep) * \real{0.3333}}
  >{\raggedright\arraybackslash}p{(\columnwidth - 4\tabcolsep) * \real{0.3333}}
  >{\raggedright\arraybackslash}p{(\columnwidth - 4\tabcolsep) * \real{0.3333}}@{}}
\toprule\noalign{}
\begin{minipage}[b]{\linewidth}\raggedright
Distinction
\end{minipage} & \begin{minipage}[b]{\linewidth}\raggedright
If conflated
\end{minipage} & \begin{minipage}[b]{\linewidth}\raggedright
Established in
\end{minipage} \\
\midrule\noalign{}
\endhead
\bottomrule\noalign{}
\endlastfoot
Admissibility (\(R\)) vs.~accessibility (\(\mathcal{R}\)) & Modal necessity collapses into a single-parameter notion that can't separately vary ``which states satisfy the regime'' from ``which states a given world can reach'' & Ch.6 \\
Inconsistency vs.~impossibility vs.~non-normality & A glutty state gets treated as automatically impossible, or an impossible state as automatically contradictory --- neither holds & Ch.7 \\
Positive support vs.~negative support & Truth values become a stipulated alphabet rather than a consequence of two independent evidential questions & Ch.9--10 \\
Support vs.~constructibility & Non-constructive derivations and exhibitable witnesses get treated as the same accomplishment & Ch.13 \\
Degree (many-valued) vs.~evidential polarity (FDE) & Łukasiewicz-style values get read as a finer-grained version of Belnap-Dunn's four, erasing that they track different things (a truth continuum vs.~two-axis evidential structure) & Ch.14 \\
Compatibility obstruction vs.~dispersion obstruction & Local disagreement at one resolution gets confused with locally-negligible variation accumulating across resolutions --- different mechanisms, different mathematics (transitivity forbids the naive version of the second) & Ch.15, Ch.20--21 \\
Restriction vs.~gluing vs.~repair & Combining a non-matching family (repair) gets described as if it were reconstructing an already-matching one (gluing) & Ch.20--21 \\
Coherence defect (\(\delta^\pm\)) vs.~glut (\(\Sigma(p)=B\)) & ``Glut = nonzero obstruction class'' gets asserted as a theorem it isn't --- proven independent instead (Mutual Non-Determination Proposition) & Ch.20--22 \\
Coherence liability vs.~coherence defect & An action's cost from evidential mismatch gets treated as a fixed function of the mismatch alone, rather than depending on what the action needs from the evidence & Ch.23 \\
Identity of states vs.~continuity of histories & Two indistinguishable present states get treated as having the same continuation possibilities, ignoring that their histories may differ & Ch.16, Ch.19 \\
State-indexed regime vs.~one global regime & Non-normal worlds get explained as declining one transformation (refusal) rather than as suspending the evaluation rule that licenses transformations generally & Ch.7 \\
\end{longtable}

\hypertarget{conclusion}{%
\chapter*{Conclusion}
\addcontentsline{toc}{chapter}{Conclusion}
\markboth{Conclusion}{Conclusion}\label{conclusion}}

\hypertarget{the-negative-results-kept-prominent-rather-than-minimized}{%
\section{The negative results, kept prominent rather than minimized}\label{the-negative-results-kept-prominent-rather-than-minimized}}

Three correspondences this project once found plausible enough to
provisionally accept did not survive their chapters, at three different
levels of the book's structure:

Three correspondences this project once found plausible enough to
provisionally accept did not survive their chapters, at three different
levels of the book's structure. At the modal level, non-normal worlds
are not refusal (Ch.7): refusal is a targeted, constitutive decision to
decline one transformation, while a non-normal world suspends the
evaluation rule itself. At the operational level, bilattice merge is not
sheaf gluing (Ch.20--21): gluing reconstructs a section from a family
that already agrees, while merge combines a family that doesn't,
retaining rather than discarding the disagreement. At the geometric
level, glut is not a nonzero cohomology class (Ch.20--22): contradiction
is a semantic fact about support, while the relevant obstruction
machinery, once built correctly, is a fact about whether evidence coheres
across overlapping contexts --- proven independent of contradiction, not
merely left unidentified with it.

These are evidence for the book's thesis, not embarrassments incidental
to it. A project that had found ways to confirm every correspondence it
initially proposed would have shown that its method could rationalize
anything; a project whose method occasionally says no is one whose
positive results --- the Mutual Non-Determination Proposition foremost
among them --- are worth taking seriously precisely because the same
method also produced these refusals.

\begin{figure}[htbp]
\centering
\begin{tikzpicture}[
  font=\small,
  ax/.style={-{Stealth[length=2.4mm]}, thick},
  coord/.style={font=\small},
  logic/.style={font=\footnotesize\itshape, gray!70!black}
]
\node[draw, circle, minimum size=13mm, inner sep=1pt] (L) at (0,0) {$\mathfrak{L}$};
\draw[ax] (L) -- (90:2.6)   node[coord, above] {$R$ \ (admissibility)};
\draw[ax] (L) -- (38:2.8)   node[coord, above right] {$\mathcal{R}$ \ (accessibility)};
\draw[ax] (L) -- (-13:3.0)  node[coord, right] {$\sigma^+, \sigma^-$ \ (support)};
\draw[ax] (L) -- (-64:2.6)  node[coord, below right] {$\kappa$ \ (constructibility)};
\draw[ax] (L) -- (-116:2.6) node[coord, below left] {$\rho$ \ (persistence)};
\draw[ax] (L) -- (-167:3.0) node[coord, left] {$\tau$ \ (continuation)};
\node[logic] at (1.3,2.35)  {non-normality varies $R_w$};
\node[logic] at (3.6,0.75)  {modal systems vary $\mathcal{R}$};
\node[logic] at (3.4,-1.15) {FDE varies $(\sigma^+,\sigma^-)$};
\node[logic] at (0.6,-2.9)  {intuitionism varies $\kappa$};
\node[logic] at (-2.7,-2.5) {vagueness varies $\rho$};
\node[logic] at (-3.5,0.6)  {the open problem lives on $\tau$};
\end{tikzpicture}
\caption{A schematic coordinate picture --- explicitly not a formal
definition --- of a logical regime
$\mathfrak{L} = (R, \mathcal{R}, \sigma^+, \sigma^-, \kappa, \rho, \tau)$
as a point in a space of independently variable structural dimensions.
The named non-classical logics do not sit along a single axis running
from classical to less classical; each alters a different coordinate,
which is this book's methodological thesis made visible. The tuple is a
mnemonic for the book's separations, not a definition the manuscript
supplies for every symbol.}
\label{fig:regime-coordinates}
\end{figure}

\hypertarget{the-open-problem-this-book-ends-on}{%
\section{The open problem this book ends on}\label{the-open-problem-this-book-ends-on}}

This is not a loose end awaiting a known fix. It's the point at which the
book's own framework ceases to close. The introduction's parameter arc
--- \(R \to (R,\mathcal{R}) \to (R,\mathcal{R},\sigma^+,\sigma^-) \to \text{local support/coherence} \to \text{repair and persistence} \to R_w\text{-dependent continuation}\) --- traced \(R\) progressively unpacking
into structure as each part of the book found the previous parameter
insufficient. Completing Parts II and V together shows what happens once
that unpacking reaches \(R\) itself: the thing Chapter 1 treated as fixed
background turns out, once identity across time is at stake, to be
capable of varying along the very trajectory it's supposed to govern.

\[x_t \xrightarrow{?} x_{t+1}, \qquad R_{x_t} \neq R_{x_{t+1}}.\]

Once admissibility is allowed to vary by state --- Chapter 7's \(R_w\),
needed to correctly place non-normal worlds --- trajectory-based identity
(Chapter 16) needs more than a sequence of admissible states connected by
\(\leadsto\). It needs an account of what makes a \emph{transition between
regimes} admissible, not merely what makes each endpoint admissible on
its own. Chapter 17's transition-object --- treating a transition as its
own kind of entity rather than a state satisfying an enriched valuation
--- is a candidate answer, floated at exactly the point in the book where
it became relevant and deliberately not promoted to a definition before
it had earned that status. Whether the right structure is a fixed choice
between \(R_{x_t}\) and \(R_{x_{t+1}}\), some combination of the two, or a
genuinely separate transition regime \(R_\gamma\), is left open. This
book's contribution is producing the question in a form precise enough
to be wrong about, not answering it --- and producing it as the place
where the framework's own regress of parameters runs out, rather than as
an item on a future-work list unconnected to everything before it.

\hypertarget{what-comes-next}{%
\section{What comes next}\label{what-comes-next}}

The correspondence ledger, not this conclusion, is the working document
for anyone extending this material --- it records status after drafting,
chapter by chapter, and should be checked before any claim in this
introduction or conclusion is relied upon for a further paper. Two
results look most likely to warrant standalone treatment ahead of the
full manuscript: the Mutual Non-Determination Proposition with its
realizability construction, and the \(R_w\)/continuation problem stated
above. Both are self-contained enough to be evaluated without the rest
of the book's apparatus, and both are the kind of result this project
was built to produce --- not agreement with Priest, and not a
redescription of him, but a specific, checkable place where this book's
own machinery says something Priest's framework alone does not.
